
在本章中，我们将证明：对于一个\trNondegenerate{}的、不依赖于时间且充分小（在 $\cC^{2}$ \trTopology{}的意义下）的\trHamiltonian{}量，当我们能够定义\trMorseComplex{}与\trFloerComplex{}时，这两者相一致。

\section{诸结果}
\label{sec:floer-to-morse-results}

出于简单性的考虑，也为了减轻书写负担，我们在第 3 章中对\trMorseComplex、而在第 9 章中对\trFloerComplex{}或多或少地使用了同一套记号。此刻，我们需要把这两者区分开来（为的是证明它们相一致）。从现在起，以 $CF_{\star}(H,J)$ 表示与\trHamiltonian{}量 $H$ 及\trAlmostComplex{}结构 $J$ 相联系的\trFloerComplex（第 9 章中记作 $(C_{\star},\partial)$），并以 $CM_{\star}(H,J)$ 表示与\trMorseFunction{} $H$ 及\trVectorField{} $\operatorname{grad}H$（即 $H$ 对于由 $J$（与 $\omega$）所定义的\trMetric{}的\trGradient）相联系的\trMorseComplex（第 3 章中记作 $(C_{\star}(H),\partial_{\operatorname{grad}H})$）。

我们想要证明下面这个结果。

\begin{theorem}\label{thm:floer-equals-morse-complex}
存在一个\trNondegenerate、充分小（就 $\cC^{2}$ \trTopology{}而言）的\trHamiltonian{}量 $H$，使得
\[ CF_{\star}(H,J)=CM_{\star+n}(H,J)\text{。} \]
\end{theorem}

这条陈述当然也包含了所论的两个\trComplex{}都是\trWellDefined{}的这一事实。从一个\trNondegenerate{}的\trHamiltonian{}量 $H_{0}$ 出发，我们把 $H$ 构造成 $H=H_{0}/k$，其中 $k$ 充分大。

由于 $H$ 充分小，它的诸\trPeriodic\trTrajectory{}都是\trStationary{}的（\cref{thm:small-hamiltonian-periodic-solutions}），于是 $\operatorname{Crit}(\scr{A}_{H})=\operatorname{Crit}(H)$。另一方面，\cref{thm:critical-nondegenerate-transfer}（它把\trNondegenerate{}性的两种概念——函数 $H$ 的诸\trCriticalPoint{}的\trNondegenerate{}性与\trHamiltonian{}量 $H$ 的诸\trPeriodic\trTrajectory{}的\trNondegenerate{}性——加以比较）蕴含：$H$ 是一个\trMorseFunction，并且 $\operatorname{Crit}(\scr{A}_{H})=\operatorname{Crit}(H)$。

由\cref{thm:linearized-flow-matrix}，我们知道 $H$ 的\trHessian{}在 $2\pi\ZZ$ 中不具有任何\trEigenvalue，于是\cref{thm:maslov-index-integer}使我们能够把 $H$ 的一个\trCriticalPoint{}的（\trMorse）\trIndex{}与相应的\trStationary\trTrajectory{}的（\trMaslov）\trIndex{}加以比较：
\[ \operatorname{Ind}_{H}(x)=\mu(x)+n\text{。} \]

这就蕴含：出现在两个\trComplex{}的定义之中的那些\trVectorSpace{}，除开\trIndex{}的一个平移之外是相同的。

接下来，我们须要考虑两个\trDifferential。更确切地说，我们必须：
\begin{itemize}
\item 能够定义这两个\trComplex{}的\trDifferential；
\item 证明它们相一致。
\end{itemize}

为了定义\trMorseComplex{}的\trDifferential，我们需要一个\trAdapted{}于\trMorseFunction{} $H$ 的\trVectorField{} $X$，它满足\trSmaleCondition。另一方面，我们又想要在这个\trVectorField{}的诸\trTrajectory{}与\trFloerEquation{}的诸解之间建立一种关系，也就是说，在
\[ \frac{du}{ds}+X(u)=0 \qquad\text{与}\qquad \frac{\partial u}{\partial s}+J(u)\frac{\partial u}{\partial t}+\operatorname{grad}H(u)=0 \]
的诸解之间建立一种关系。因此，我们想要让 $X$ 是函数 $H$ 对于由某个受 $\omega$ \trCalibrated{}的\trAlmostComplex{}结构 $J$ 所定义的\trMetric{}的\trGradient。于是我们将证明下面这个结果。

\begin{theorem}\label{thm:calibrated-dense-morse-smale}
设 $H$ 是\trSymplectic\trManifold{} $(W,\omega)$ 上的一个\trMorseFunction。那么存在一个由受 $\omega$ \trCalibrated{}的\trAlmostComplex{}结构所组成的\trDense{}子集 $\scr{J}_{\mathrm{reg}}(H)$，使得偶对 $(H,-JX_{H})$ 是一个\trMorse--\trSmale{}对。
\end{theorem}

我们将分两步进行证明。第一步将致力于一个一般的\trMorseFunction{} $f$ 与一个\trAdapted{}的\trVectorField{} $X$。我们将把\trVectorField{} $-X$ 的\trFlow{}的方程
\[ \frac{du}{ds}+X(u)=0 \]
沿它的一个解\trLinearization{}为 $L_{u}Y=0$，然后证明下面这个结果。

\begin{theorem}\label{thm:lu-fredholm-index-difference}
当 $f$ 是一个\trMorseFunction{}且 $u$ 是一条连接两个\trCriticalPoint{} $x$ 与 $y$ 的\trTrajectory{}时，$L_{u}$ 是一个\trFredholmOperator，其\trIndex{}是这两个\trCriticalPoint{}的\trMorse{}\trIndex{}之差。
\end{theorem}

由于前面提到过的\trMorse{}\trIndex{}与\trMaslov{}\trIndex{}之间的关系，这就立即蕴含下面这条推论。

\begin{corollary}\label{thm:df-lu-same-index}
对于一个\trNondegenerate{}的\trHamiltonian{}量 $H$ 与 $-JX_{H}$ 的一条\trTrajectory{} $u$，诸\trFredholmOperator{} $(d\scr{F})_{u}$ 与 $L_{u}$ 具有相同的\trIndex。
\end{corollary}

我们还将证明下面这个结果。

\begin{theorem}\label{thm:smale-iff-lu-surjective}
\trVectorField{} $X$ 满足\trSmaleCondition，当且仅当诸算子 $L_{u}$ 都是\trSurjective{}的。
\end{theorem}

在第二步中，我们将考虑\trHamiltonian{}量 $H$ 这一具体情形，并证明\cref{thm:calibrated-dense-morse-smale}——这是一条\trTransversality{}结果，与我们曾在\cref{sec:transversality}中建立的那些结果相类似。在选定了\trAlmostComplex{}结构（从而也选定了\trSmale{}\trVectorField）之后，我们随后就将把\trFloerEquation{}的诸解与这个\trVectorField{}的诸\trTrajectory{}加以比较。

\begin{remark}\label{thm:t-independent-floer-solutions}
\trFloerEquation{}的那些不依赖于 $t$ 的解，恰好就是\trVectorField{} $X=-\operatorname{grad}H$ 的诸\trTrajectory，这一点是显然的。
\end{remark}

我们想要证明下面这条命题。

\begin{proposition}\label{thm:kernel-df-equals-lu}
若 $H$ 充分小（在 $\cC^{2}$ 意义下），那么
\[ \operatorname{Ker}(d\scr{F})_{u}=\operatorname{Ker}L_{u}\text{。} \]
\end{proposition}

我们将在后面证明这一点。于是，$(d\scr{F})_{u}$ 的\trKernel{}的元素都不依赖于 $t$。利用诸\trIndex{}的同一性（\cref{thm:df-lu-same-index}以及\cref{thm:smale-iff-lu-surjective}中对\trSmale{}\trVectorField{}的刻画），这就蕴含下面这个结果。

\begin{corollary}\label{thm:df-surjective-along-gradient}
\trFredholmOperator{} $(d\scr{F})_{u}$ 沿 $H$ 的\trGradient{}的每一条\trTrajectory{}都是\trSurjective{}的。
\end{corollary}

另一方面，我们将证明：把 $H$ 换成 $H_{k}=H/k$（它具有相同的\trCriticalPoint{}与相同的\trIndex）之后——如有必要，取 $k$ 充分大——我们所关心的\trFloerEquation{}的一切解都不依赖于 $t$。

\begin{proposition}\label{thm:hk-large-t-independent}
若 $k$ 充分大，那么 $H_{k}$ 的\trFloerEquation{}的、连接满足
\[ \operatorname{Ind}_{H_{k}}(x)-\operatorname{Ind}_{H_{k}}(y)\le 2 \]
的\trCriticalPoint{} $x$ 与 $y$ 的那些解，全都不依赖于 $t$。
\end{proposition}

设 $H_{k}=H/k$，其中 $k$ 充分大，并设 $J\in\scr{J}_{\mathrm{reg}}$（由\cref{thm:calibrated-dense-morse-smale}所给出的那个\trDense{}子集）。结论便是：当 $\operatorname{Ind}(x)-\operatorname{Ind}(y)\le 2$ 时，与 $(H,J)$ 相联系的\trFloerEquation{}的、连接两个\trCriticalPoint{} $x$ 与 $y$ 的那些\trTrajectory（由于两个\trCriticalPoint{}集相重合，它们既是 $H$ 的\trCriticalPoint{}也是 $\scr{A}_{H}$ 的\trCriticalPoint）恰好就是\trSmale{}\trVectorField{} $X=-JX_{H}$ 的诸\trTrajectory（由\cref{thm:hk-large-t-independent}）。此外，沿这些\trTrajectory{}\trLinearization{}的\trFloer{}算子是\trSurjective{}的（由\cref{thm:df-surjective-along-gradient}）。这样一条\trRegularity{}条件蕴含 $\scr{M}^{(H,J)}(x,y)$ 是一个\trManifold，这就使我们能够定义\trFloerComplex{}（参见\cref{thm:moduli-manifold-dimension}）。

我们只对 $\operatorname{Ind}(x)-\operatorname{Ind}(y)\le 2$ 建立这条\trRegularity{}性质（而不是像第 8 章中那样对一般情形建立，参见\cref{thm:regular-perturbations}），但这一情形已足以构造\trFloerComplex{} $CF_{\star}(H,J)$。我们已经看到，它与\trMorseComplex{} $CM_{\star+n}(H,J)$ 相一致，这就证明了\cref{thm:floer-equals-morse-complex}。我们还须证明那些中间结果（\cref{thm:calibrated-dense-morse-smale}、\cref{thm:lu-fredholm-index-difference}、\cref{thm:smale-iff-lu-surjective}与\cref{thm:kernel-df-equals-lu}）。

